% The axis the upkeep term rests on, and the four refining axes. Inline
% in the preprint; the venue build keeps the two deciding axes and the
% summary in the body and moves these here.
\paragraph{Upkeep is not a modelling convenience.}
The premise under all of this is that memory an agent carries but does not
earn from has a real, recurring cost. \citet{gloaguen2026agentsmd} price it
directly on the surface where agent memory is most widely deployed:
repository-level context files, always loaded, never curated. Across four
models and two task sets they report that context files increase inference
cost by over 20\% on average while not generally improving task success, and
that repository overviews---the section every provider's initialisation prompt
recommends---are not helpful. Two things about that measurement bound what we
may take from it. It compares against \emph{no context file}, and
developer-written files do beat model-written ones, so it is not a finding
that written-down memory is worthless. And it prices a file injected into
every prompt, not an entry fetched by retrieval, so it is a strong analogy for
our upkeep term rather than a measurement of it. What it does establish is
that the cost side of the ledger is not a modelling device we introduced to
make selection work: somebody else measured it, on a mechanism with no
curation at all, and found it substantial.


\paragraph{Axis 2a: If measurements lie on purpose, \emph{what is the
liar buying}?} Our adversary was written to buy destruction: it spends
its budget on every measured outcome, blaming benign entries as well as
paying the poison. A deployed system suggested the other objective is
the one an attacker would actually choose (\S\ref{sec:memoryos}), so we
ran it --- an adversary identical in every respect except that it lies
\emph{only} when the poison has just done damage.

Persistence is a real attack on the ledger and we report it as one. Its
poison-kill rate falls from $1.00$ to $0.10$ at two lies per cycle and
to $0.00$ at four (paired $p = 0.0039$, 10 seeds), and it does so
without moving benign capability at all, so nothing appears in the
metric an operator watches. The cause is our own central mechanism read
from the other side: bounded credit with earn-back is what forgives a
lie the defender did not deserve, and equally what lets a \emph{paid}
poison earn its way back above the floor.

What it does \emph{not} do is reverse the recommendation, and getting
that wrong is a trap this paper has already fallen into once
(\S\ref{sec:swebench-attack}). The strike counters' kill rate is
unmoved by the attack, which reads as immunity until one counts poison
by \emph{provenance} rather than by whether it currently advises action.
On that predicate the counters never removed the poison at all: two
poisoned entries survive in \texttt{evict\_on\_negative} at every
budget \emph{including zero}, against none for the unattacked ledger,
and quarantine ends worse under attack than it began ($2.0
\rightarrow 2.8$). The ledger under the heaviest attack we run still
holds fewer poisoned entries ($1.0$) than any counter holds with no
attacker present ($2.0$; mean difference $-1.0$, $p = 0.0020$). A
mechanism that keeps the inert poison forever has nothing left for a
persistence adversary to take, which is abstention rather than defence
--- the same reading we give the bandit and \texttt{keep\_everything}
on Axis 2.

The honest statement is therefore narrower than a flip: \emph{persistence
narrows the ledger's advantage from total elimination to partial, and
costs it the guarantee, without making any counter preferable.} The
attack is worth publishing because a defence whose guarantee can be
removed silently is one an operator should know about, not because the
alternative is better.

\paragraph{Axis 3: Do measurements lie by accident?}
When measurements are truthful and deterministic, a one-line
``evict-on-negative'' heuristic \emph{ties} the full ledger on outcomes
(Table~\ref{tab:headline}); the ledger's extra machinery buys only
leanness there. As measurements become noisy, the picture splits by
\emph{shape}. Under one-sided false-bad noise (good actions sometimes
misreport as bad---the flaky-CI case), the ledger's bounded-credit buffer
with earn-back is decisively better than any fixed strike counter
(Table~\ref{tab:noise}), and matches an optimistic confidence bandit only
at high rates. Under symmetric noise that also rewards the guilty, every
mechanism has a failure boundary; we locate the ledger's at roughly one
lie in three and publish it. The honest rule: \emph{if your measurements
never lie and you do not need leanness, the if-statement is the right
tool}; the ledger earns its complexity exactly when measurements are
noisy in a one-sided way, which is the common CI failure mode.

\paragraph{Axis 4: Are refusals measured?}
This axis decides the second-environment reversal. The ledger's leanness
is an asset when inaction produces a measurable outcome (storage: keeping
a file is measured) and a liability when refusing to act is free
(test-suite: a skipped patch runs no suite, so a protective entry that
keeps refusing a destructive change can never earn and starves). Where
refusals are unpriced, a mechanism that never removes entries (a strike
counter, or the bandit) outlasts the ledger because its protector refuses
forever. The design implication is concrete: pair outcome-based
selection with an environment that prices inaction (e.g.\ a
reverted-incident counter that pays a blocked bad change), or with a
quarantine-style cooldown, and the reversal is expected to close. We do
not claim it; the environment that would show it is future work.

\paragraph{Axis 5: Is the horizon longer than the time constant?}
The mechanism has a clock, and it is the first thing to check. An entry
spawns with $1.0$ energy and pays $0.05$ upkeep, so removal by disuse
takes $20$ cycles; nothing the ledger does that a simpler policy does not
can appear before then. Our own real-task leg violated this and we did
not notice until it returned a null: SWE-Bench-CL's two pinned sequences
are $19$ and $22$ tasks, $10$ of the $246$ lessons minted by the two
upkeep-charging arms ever starved, and
\texttt{memory\_on} tied \texttt{keep\_everything} to three decimals
because on that horizon they are the same policy
(\S\ref{sec:swebench}). The rule is blunt: if the deployment is shorter
than a few multiples of $\textrm{spawn}/\textrm{upkeep}$, use the
if-statement---the ledger cannot yet have done anything else---or raise
upkeep so the time constant fits the horizon.

\paragraph{Axis 6: What do you need besides correctness?}
Three needs select the ledger regardless of the above, because no other
mechanism we tested provides them. \emph{Leanness}: only the energy
ledger starves never-consulted entries; counters and the bandit hoard
them (population 15--16 vs 4). \emph{Completeness}: among arms that
retain any capability, the ledger is the only one that ends with no
poisoned entry alive, acting or dormant---every counter, quarantine
policy and bandit we test ends holding two or three poisoned entries at
\emph{every} attack budget including zero, and \texttt{ttl} matches the
ledger's zero only by emptying the store (benign capability $0.00$). Two
scope conditions belong on this claim rather than in a footnote, because
both are measured in this paper. It holds \emph{unattacked}: a
persistence adversary takes the ledger to $1.0$ poisoned entries alive
(\S\ref{sec:persistence}), which is still fewer than any counter holds
with no attacker present, but it is not zero. And it does not hold on
real tasks: on SWE-Bench-CL no arm removed poison at all, the ledger
included (\S\ref{sec:swebench-attack}). Note also that this is
completeness and not speed---in the deterministic suites it kills
actionable poison within one cycle, but with a real model deciding
adoption a one-line counter revokes five to nine cycles sooner and the
ledger's advantage is that it finishes the job
(\S\ref{sec:wef}). \emph{Cost}: updating balances from a reported scalar is
$10^4$--$10^6\times$ cheaper than an LLM judge or generator, with no API
dependency and no determinism loss.

